% !TEX root = ../main.tex
\chapter{Backend Dispatch}
\label{ch:backend_dispatch}

A tensor framework's dispatch layer is the indirection that maps a
device-agnostic operation request --- ``add these two tensors'' ---
to a device-specific implementation. In \Lucid\ the dispatch
layer is small and explicit: a single \code{Dispatcher} class
inspects the input tensors' device tag and routes to one of two
concrete backends, with two well-defined carve-outs that route a
device-tagged tensor through the opposite backend's path. This
chapter formalises the dispatcher's contract, describes the two
backends and the carve-outs, and explains how the Python side
mirrors the C++ structure.

The dispatcher is also where the \emph{H3 hard rule} from
\appref{app:hard_rules} is enforced: CPU stream uses only \Accelerate,
GPU stream uses only \MLX, with the two named exceptions. Any new
backend code that violates the rule is rejected at review. The
discipline is what keeps the engine simple at the cost of writing
fewer one-off kernels.

The chapter proceeds: \secref{sec:backend_dispatch:ibackend}
describes the \code{IBackend} interface;
\secref{sec:backend_dispatch:dispatcher} describes the
\code{Dispatcher} class; \secref{sec:backend_dispatch:cpu_gpu}
describes the two backends; \secref{sec:backend_dispatch:carveouts}
describes the linalg and data-dependent-output carve-outs;
\secref{sec:backend_dispatch:mpsgraph_carveout} treats the 3.4 per-op
\MPSGraph\ dispatch; \secref{sec:backend_dispatch:dtype_promo}
describes type promotion; \secref{sec:backend_dispatch:python}
treats the Python side; \secref{sec:backend_dispatch:perf} closes
with the cost model.

\section{The IBackend Interface}
\label{sec:backend_dispatch:ibackend}

\code{IBackend} (in \code{lucid/\_C/backend/IBackend.h}) is the
abstract interface every backend implements. It is a large
interface --- one method per op category --- but each method has a
narrow contract:

\begin{lstlisting}[style=lucidcpp,language=C++]
class IBackend {
public:
  virtual ~IBackend() = default;

  // Memory
  virtual TensorImpl allocate(Shape s, Dtype dt) = 0;
  virtual void       free(TensorImpl& t)         = 0;

  // Elementwise
  virtual TensorImpl unary (OpKind, const TensorImpl&)               = 0;
  virtual TensorImpl binary(OpKind, const TensorImpl&, const TensorImpl&) = 0;

  // Reductions
  virtual TensorImpl reduce(OpKind, const TensorImpl&,
                            AxisSet, bool keepdims) = 0;

  // Matrix multiplication
  virtual TensorImpl matmul(const TensorImpl& a, const TensorImpl& b) = 0;

  // Convolution
  virtual TensorImpl conv2d(const TensorImpl& x, const TensorImpl& w,
                            ConvParams) = 0;
  // ... and so on, one method per operation category
};
\end{lstlisting}

The interface is intentionally op-category-granular, not per-op.
Within \code{unary}, the \code{OpKind} enum distinguishes
\code{Relu}, \code{Sigmoid}, \code{Tanh}, etc. This keeps the
interface from exploding to hundreds of methods, at the cost of an
internal switch inside each method. The switch is small and
predictable; modern CPUs predict it perfectly after the first
invocation.

\subsection{Why a virtual interface}
\label{ssec:backend_dispatch:why_virtual}

The virtual-call overhead is roughly 2--5 nanoseconds per dispatch
on modern Apple silicon. Compared to the cost of any nontrivial
operation it dispatches to (microseconds for a small GPU launch,
hundreds of microseconds for a matrix multiplication), this is
negligible.

A template-based static dispatch would have been the alternative.
We chose virtual dispatch for two reasons: (i)~the binary size
under templates is much larger (every op kind generates new code
per backend); (ii)~runtime device selection requires runtime
dispatch anyway, since the user can change a tensor's device
between calls.

\section{The Dispatcher Class}
\label{sec:backend_dispatch:dispatcher}

\code{Dispatcher} (in \code{lucid/\_C/backend/Dispatcher.h}) is
the front door. It holds pointers to the two concrete backends and
routes every operation through the appropriate one.
\figref{fig:backend_dispatch:routing} sketches the full routing
graph including the two H3 carve-outs and the 3.4 per-op
\MPSGraph\ shim.

\begin{figure}[t]
\centering
\begin{adjustbox}{max width=\textwidth, center}
\begin{tikzpicture}[
  in/.style={draw, rounded corners=2pt, font=\small\sffamily,
             fill=lucid.note.bg, minimum width=3.4cm,
             minimum height=0.9cm, align=center, inner sep=4pt},
  dec/.style={draw, diamond, aspect=2, font=\small\sffamily,
              fill=lucid.code.bg, inner sep=2pt, align=center,
              minimum width=2.6cm},
  be/.style={draw, rounded corners=2pt, font=\small\sffamily,
             fill=lucid.info.bg, minimum width=3.6cm,
             minimum height=0.9cm, align=center, inner sep=4pt},
  cv/.style={draw, rounded corners=2pt, font=\small\sffamily,
             fill=lucid.ok.bg, minimum width=3.6cm,
             minimum height=0.9cm, align=center, inner sep=4pt},
  arrow/.style={->, >=stealth, thick, color=lucid.accent.dark,
                shorten >=2pt, shorten <=2pt},
  lbl/.style={font=\footnotesize\itshape, color=lucid.muted,
              fill=white, inner sep=2pt},
]
  \node[in] (op) at (0, 6) {Op request\\(OpKind, tensors)};
  \node[dec] (dev) at (0, 4) {device?};

  \node[dec] (cv1) at (-6.5, 2) {linalg?\\data-dep?};
  \node[be] (accel) at (-6.5, -0.5) {AccelerateBackend\\(vDSP/BLAS/BNNS)};

  \node[dec] (cv2) at (6.5, 2) {MPSGraph\\carve-out?};
  \node[be] (mlx) at (6.5, -0.5) {MlxBackend\\(\MLX\ primitives)};
  \node[cv] (mpsg) at (6.5, -3) {MpsBuilder\\(per-op compile)};

  \draw[arrow] (op) -- (dev);
  \draw[arrow] (dev.west) to[bend right=10]
    node[lbl, midway, above] {cpu} (cv1.north);
  \draw[arrow] (dev.east) to[bend left=10]
    node[lbl, midway, above] {metal} (cv2.north);

  \draw[arrow] (cv1) -- node[lbl, right] {no} (accel);
  \draw[arrow] (cv1.south) to[bend right=15]
    node[lbl, midway, below] {linalg / data-dep} (mlx.north west);

  \draw[arrow] (cv2) -- node[lbl, right] {no} (mlx);
  \draw[arrow] (cv2.east) to[bend left=30]
    node[lbl, midway, right] {gelu, layernorm, \ldots} (mpsg.east);
\end{tikzpicture}
\end{adjustbox}
\caption[Dispatcher routing.]{Op-request routing inside
\Lucid's dispatcher. The primary decision is the input tensor's
\code{device} tag. The two H3 carve-outs (linalg and
data-dependent output ops) cross from the cpu branch to the metal
branch under the hood; the 3.4 per-op \MPSGraph\ carve-out
intercepts a small set of metal-device ops and routes them through
the \code{MpsBuilder} compilation path instead of \MLX.}
\label{fig:backend_dispatch:routing}
\end{figure}

\tabref{tab:backend_dispatch:carveouts} summarises the carve-outs.

\begin{table}[t]
\centering
\footnotesize
\setlength{\tabcolsep}{4pt}
\begin{adjustbox}{max width=\textwidth, center}
\begin{tabular}{lllll}
\toprule
Carve-out & Ops affected & User device & Actual path & Rationale \\
\midrule
linalg          & cholesky, qr, svd, eigh, ... & metal $\to$ cpu & MLX (CPU stream) & no Metal LAPACK \\
data-dep output & nonzero, unique, masked\_select & metal $\to$ cpu & Accelerate & shape needs value read \\
3.4 MPSGraph    & gelu, layernorm, embedding\_bwd, ... & metal & MPSGraph & 3--30$\times$ over MLX \\
\bottomrule
\end{tabular}
\end{adjustbox}
\caption[Dispatcher carve-outs.]{The three named carve-outs in
\Lucid's dispatcher. Each is justified by a structural reason
that the framework cannot work around at the dispatch level.}
\label{tab:backend_dispatch:carveouts}
\end{table}

\begin{lstlisting}[style=lucidcpp,language=C++]
class Dispatcher {
public:
  static Dispatcher& instance();

  TensorImpl unary(OpKind op, const TensorImpl& x) {
    return select_backend(x.device)->unary(op, x);
  }
  TensorImpl binary(OpKind op, const TensorImpl& a, const TensorImpl& b) {
    if (a.device != b.device) {
      throw LucidError("binary op across devices");
    }
    return select_backend(a.device)->binary(op, a, b);
  }
  // ... and so on

private:
  IBackend* select_backend(Device d) {
    return (d == Device::Cpu) ? &accelerate_backend_ : &mlx_backend_;
  }

  AccelerateBackend accelerate_backend_;
  MlxBackend        mlx_backend_;
};
\end{lstlisting}

The dispatcher is a singleton. The two backend members are
constructed eagerly at the singleton's first access; constructing
them takes a few microseconds each (mostly \MLX's initialisation).

\subsection{Cross-device errors}
\label{ssec:backend_dispatch:cross_device}

Binary and higher-arity operations require all inputs to share a
device. The dispatcher checks this at the front door and raises
\code{LucidError("binary op across devices")} on mismatch, with a
message that names the operation and the two devices.

The user-side convention is to \code{.to(device)} explicitly before
the operation, rather than to expect implicit moves. \Lucid\ does
not implicitly move tensors across devices for two reasons: the
move has a real cost (the bridge of \chref{ch:unified_memory} is
free only in the contiguous case), and the move's direction is not
unambiguous (which device should the result land on?).

\subsection{Per-op routing inside the backend}
\label{ssec:backend_dispatch:per_op_internal}

Once the dispatcher hands a request to a backend, the backend uses
the \code{OpKind} enum to switch to the per-op implementation. For
\code{AccelerateBackend}, the switch table routes to vDSP / vForce
/ BLAS / BNNS helpers in \code{AccelerateMath.cpp} and friends. For
\code{MlxBackend}, the switch routes to \MLX\ primitive calls in
\code{MlxBackend.cpp}.

The switch is one indirection layer that could in principle be
collapsed (a backend could expose one method per op) but is kept
for code organisation: a 260-op backend would have 260 virtual
methods, which is unmanageable.

\section{The Two Backends}
\label{sec:backend_dispatch:cpu_gpu}

The two concrete backends are described in detail in the hardware
chapters: \chref{ch:accelerate} for the CPU side and
\chref{ch:mlx_runtime} for the GPU side. This section gives the
short version from the dispatcher's perspective.

\subsection{AccelerateBackend}
\label{ssec:backend_dispatch:accelerate}

\code{AccelerateBackend} is the CPU implementation of
\code{IBackend}. Every method routes to an \Accelerate-backed
helper:

\begin{itemize}
  \item Elementwise unary/binary $\to$ vDSP primitives (for
    arithmetic) or vForce primitives (for transcendentals).
  \item Reductions $\to$ vDSP reductions (\code{vDSP\_sve},
    \code{vDSP\_meanv}, etc.).
  \item \code{matmul} $\to$ \code{cblas\_sgemm} or
    \code{cblas\_dgemm}.
  \item \code{conv2d} $\to$ BNNS direct convolution.
  \item Allocation $\to$ \Lucid's pool wrapping the system
    allocator.
\end{itemize}

\code{AccelerateBackend} does not link against \MLX. It is the
self-contained CPU path.

\subsection{MlxBackend}
\label{ssec:backend_dispatch:mlx}

\code{MlxBackend} is the GPU implementation. Every method routes
to an \MLX\ primitive call with \code{stream=mx.gpu}:

\begin{itemize}
  \item Elementwise $\to$ \code{mlx::core::add}, etc.
  \item Reductions $\to$ \code{mlx::core::sum}, etc.
  \item \code{matmul} $\to$ \code{mlx::core::matmul}.
  \item \code{conv2d} $\to$ \code{mlx::core::conv2d}.
  \item Allocation $\to$ \Lucid's pool wrapping
    \code{mlx::core::array}'s storage.
\end{itemize}

\code{MlxBackend} does not link against \Accelerate\ for its own
operations, except indirectly through \MLX\ itself (which uses
\Accelerate\ in its CPU stream).

\section{The Two Carve-Outs}
\label{sec:backend_dispatch:carveouts}

The H3 rule (\appref{app:hard_rules}) has two named carve-outs.
Both are documented here.

\subsection{The linalg carve-out}
\label{ssec:backend_dispatch:linalg_carveout}

Linear-algebra decompositions (\code{cholesky}, \code{qr},
\code{svd}, \code{eigh}, \code{solve}, \code{lstsq}, \code{inv},
\code{det}, \code{logdet}, \code{matrix\_exp}) are not implemented
in \Metal. \MLX\ provides them only on its CPU stream, where they
delegate to \Accelerate's LAPACK. \Lucid's dispatcher recognises
this and routes linalg ops through \MLX-on-CPU regardless of the
user's chosen device:

\begin{lstlisting}[style=lucidcpp,language=C++]
TensorImpl Dispatcher::cholesky(const TensorImpl& a) {
  // Linalg carve-out: always CPU stream, even for metal tensors.
  TensorImpl a_cpu = (a.device == Device::Metal)
      ? bridge_metal_to_cpu(a)
      : a;
  TensorImpl r_cpu = mlx_backend_.cholesky(a_cpu, /*stream=*/mx::cpu);
  return (a.device == Device::Metal)
      ? bridge_cpu_to_metal(r_cpu)
      : r_cpu;
}
\end{lstlisting}

The bridge is the zero-copy \code{SharedStorage} of
\chref{ch:unified_memory}, so the round-trip is free in the
contiguous case. The carve-out's user-visible effect is that
\code{lucid.linalg.cholesky(x)} on a metal tensor takes the same
time it would on a CPU tensor, plus a negligible bridge overhead;
the device argument is effectively fictive for these operations.

The chapter that fills in the linalg subsystem is
\chref{ch:linalg}.

\subsection{The data-dependent output carve-out}
\label{ssec:backend_dispatch:data_dep_carveout}

A second class of operations produces output whose shape depends on
the input \emph{values}, not just the input shape. The canonical
examples are \code{nonzero}, \code{unique}, and
\code{boolean\_mask\_select}. \MLX\ supports these on the GPU
stream, but the runtime cost includes a synchronous read of the
input values to determine the output shape, which forces a
GPU-to-CPU stall. The stall is unavoidable in any framework on any
platform; it is fundamental to the operation.

\Lucid's dispatcher recognises this and explicitly bridges to the
CPU for the value read, then back to the GPU for the output:

\begin{lstlisting}[style=lucidcpp,language=C++]
TensorImpl Dispatcher::nonzero(const TensorImpl& x) {
  if (x.device == Device::Metal) {
    TensorImpl x_cpu = bridge_metal_to_cpu(x);
    TensorImpl r_cpu = accelerate_backend_.nonzero(x_cpu);
    return bridge_cpu_to_metal(r_cpu);
  }
  return accelerate_backend_.nonzero(x);
}
\end{lstlisting}

The user-visible effect is that data-dependent ops on metal tensors
are not faster than on CPU tensors. We document this in the API
reference (\chref{ch:indexing}); the small cost is the price of
supporting the operation at all.

\section{The Per-Op MPSGraph Carve-out}
\label{sec:backend_dispatch:mpsgraph_carveout}

\Lucid\ 3.4 introduced a third carve-out: a small set of GPU
operations dispatch through \MPSGraph\ instead of \MLX, because the
\MPSGraph\ implementations are substantially faster for those ops.
The current set:

\begin{itemize}
  \item GELU (forward and backward).
  \item LayerNorm (forward and backward).
  \item Embedding backward.
  \item MaxPool backward.
  \item SiLU (forward and backward).
  \item SoftmaxCrossEntropyWithLogits.
\end{itemize}

The dispatcher's GPU path checks for these ops first and, if
matched, constructs a one-op \MPSGraph, compiles it, and runs it
through the MpsBuilder infrastructure
(\chref{ch:mps_builder}). The compiled graph is cached so that
subsequent invocations with the same input shape skip recompilation.

The carve-out is invisible to the user: \code{nn.GELU} produces the
same numerical result regardless of whether it dispatched through
\MLX\ or \MPSGraph. The chapter that documents this in detail is
\chref{ch:mpsgraph_dispatch}.

\subsection{When to add a new MPSGraph carve-out}
\label{ssec:backend_dispatch:add_mpsgraph}

The policy is to add a new op to the per-op \MPSGraph\ carve-out
when the empirical speedup over the \MLX\ implementation exceeds a
factor of three on a representative training workload and when the
\MPSGraph\ implementation matches \MLX's numerical behaviour to
within $10^{-5}$ relative tolerance for the dtype.

Both criteria are necessary: a small speedup does not justify the
maintenance overhead of a parallel implementation, and numerical
divergence would silently change training trajectories.

\section{Type Promotion}
\label{sec:backend_dispatch:dtype_promo}

The dispatcher also handles dtype promotion: when a binary operation
receives inputs of different dtypes, the dispatcher promotes both
to a common dtype before calling the backend. The promotion rules
mirror NumPy's:

\begin{itemize}
  \item Floats win over ints (\code{float32 + int32 $\to$ float32}).
  \item Wider wins over narrower
    (\code{float16 + float32 $\to$ float32}).
  \item Signed wins over unsigned at equal width (rare in \Lucid,
    which does not expose unsigned integers, but the rule is
    preserved for \MLX-side inter-op).
  \item \code{bfloat16 + float16 $\to$ float32} (the two narrow
    floats are not directly compatible; promotion to FP32 is the
    only safe choice).
\end{itemize}

The rules are encoded in
\code{lucid/\_ops/composite/dtype.py} as a promotion table; the
dispatcher consults the table at every binary op call. The table is
small enough to fit in L1 cache.

\subsection{Promotion in AMP mode}
\label{ssec:backend_dispatch:amp_promo}

Mixed precision (AMP, \chref{ch:amp}) overrides the default
promotion in two ways: under \code{autocast}, certain operations
are forced to FP16 even when their inputs are FP32, and certain
operations (typically reductions) are forced to FP32 even when
inputs are FP16. The override is implemented as an AMP-aware shim
in the dispatcher; we treat the shim's details in
\chref{ch:amp_integration}.

\section{The Python Side}
\label{sec:backend_dispatch:python}

The C++ dispatcher described above is reached through a Python-side
dispatcher in \code{lucid/\_dispatch.py}, which performs the
\code{\_wrap} / \code{\_unwrap} between Python \code{Tensor} and
engine \code{TensorImpl}.

\subsection{The wrap/unwrap pair}
\label{ssec:backend_dispatch:wrap_unwrap}

The two functions:

\begin{lstlisting}[style=lucid,language=LucidPython]
from lucid._C import engine as _C_engine

def _unwrap(obj):
    """Tensor -> TensorImpl, or scalar -> 0d TensorImpl."""
    if isinstance(obj, Tensor):
        return obj._impl
    if isinstance(obj, (int, float, bool, complex)):
        return _C_engine.scalar_tensor(obj)
    raise TypeError(...)

def _wrap(impl):
    """TensorImpl -> Tensor."""
    t = Tensor.__new__(Tensor)
    t._impl = impl
    return t
\end{lstlisting}

Note the import: \code{from lucid.\_C import engine as \_C\_engine}.
The alias is mandatory per Hard Rule H2 (\appref{app:hard_rules});
bare \code{import engine} or \code{as e} is rejected at review.

\subsection{Per-op wrappers}
\label{ssec:backend_dispatch:py_wrappers}

Each Python-side op wrapper is a thin function that unwraps inputs,
calls the C++ dispatcher (which is itself bound through pybind11),
and wraps the result:

\begin{lstlisting}[style=lucid,language=LucidPython]
def add(a, b):
    return _wrap(_C_engine.add(_unwrap(a), _unwrap(b)))
\end{lstlisting}

The wrappers are auto-generated for elementwise ops (a single
template fills in 60+ unary/binary functions); only ops with
nontrivial Python-side adaptation (e.g.\ argument-name remapping,
default-value handling) are written by hand.

The result is that the Python overhead per op call is a few
microseconds: argument validation, unwrap, C++ call, wrap. This is
dominated by the cost of any nontrivial operation but matters for
elementwise ops on small tensors.

\section{Performance Characteristics}
\label{sec:backend_dispatch:perf}

The end-to-end cost of a single op call, in nanoseconds on M3~Max:

\begin{itemize}
  \item Python wrapper (unwrap, argument validation): $\approx 800$
    ns.
  \item Pybind11 boundary (Python $\to$ C++): $\approx 200$ ns.
  \item C++ dispatcher (device check, switch to backend): $\approx
    20$ ns.
  \item Backend per-op switch (\code{OpKind} to implementation):
    $\approx 5$ ns.
  \item Provider call (\MLX\ primitive or vDSP function): varies;
    typically 200\,ns to several $\mu$s.
  \item Pybind11 boundary back: $\approx 200$ ns.
  \item Python wrap: $\approx 200$ ns.
\end{itemize}

Total framework overhead per op call: roughly 1.5\,$\mu$s, of which
half is Python-side. For elementwise ops on small tensors, the
overhead can match the kernel cost; for medium and large tensors,
the overhead is negligible.

\Lucid\ 3.5's compile path (\chref{ch:compile_design}) eliminates
the per-op overhead inside compiled regions: the dispatcher runs
once during tracing and the compiled graph then executes without
per-op Python/C++ boundary crossings. The win is most visible on
small-tensor-heavy workloads where the overhead would otherwise
dominate.

\section{Summary and Forward References}
\label{sec:backend_dispatch:summary}

The backend dispatch layer in \Lucid\ is small and explicit: a
single \code{Dispatcher} routes per-op requests to one of two
\code{IBackend} implementations (\code{AccelerateBackend} for CPU,
\code{MlxBackend} for GPU), with two named carve-outs (linalg
routes through CPU stream regardless of device; data-dependent
output ops also route through CPU) and a third carve-out for the
3.4 per-op \MPSGraph\ dispatch. The Python side mirrors the C++
structure through \code{lucid/\_dispatch.py}, with the H2-required
\code{\_C\_engine} alias.

The next chapter, \chref{ch:bridge_boundaries}, catalogues the six
places where \Lucid\ touches the world outside the engine.
\chref{ch:python_cpp_binding} then describes the pybind11 layer
that the dispatcher sits on top of. \chref{ch:op_registry} returns
to the dispatch question from the op-registration angle.

\begin{lucidnote}
For the user: device choice (\code{cpu} vs \code{metal}) controls
which backend the dispatcher selects, with the linalg carve-out
making the choice fictive for decomposition ops. Cross-device
binary operations raise an error rather than implicit-moving; the
user is expected to \code{.to(device)} explicitly.
\end{lucidnote}
